\chapter{Cycle-Overlap Rigidity}

\section{Cycle Space}

Let \(G=(V,E)\) be a finite undirected graph and let \(c(G)\)
be its number of connected components. Its cycle-space dimension is
\[
\beta_1(G)=|E|-|V|+c(G).
\]
For a connected graph this reduces to \( |E|-|V|+1 \).
Indeed, a spanning forest has \( |V|-c(G) \) edges, and the
fundamental cycles associated with the remaining edges form a basis.

\section{Cycle Incidence Matrix}

Define \(A\) where rows correspond to cycles.

\section{Rigidity Statement: Corrected Scope}

The former unrestricted assertion that increasing cycle-overlap rank
forces failure of FO$^k$-local homogeneity is not established.
For the interpretation as the span of embedded radius-two cycle spaces,
it is false. For even \(m\ge6\), the toroidal grid
\(C_m\square C_m\) is vertex-transitive and has maximum degree four.
The \(m^2/4\) unit squares with lower-left coordinates \((2i,2j)\)
have pairwise disjoint vertex sets and independent edge-boundary vectors.
Each lies in a radius-two ball. Thus the local cycle-span rank grows
without bound while every rooted logical type remains identical.

A bound on the rank of a single fixed-radius ball follows from the
degree bound alone. Any stronger rigidity claim requires a precise
overlap invariant and additional hypotheses.

